\section{战国纲要：各国都在学会变强}

战国不是“七个国家打来打去”这么简单。各国都发现，春秋式的封邑、卿族和临时会盟不足以支撑长期竞争。要活下去，国家得知道有多少人、能收多少粮、能征多少兵、怎样奖功罚罪、怎样让命令走到县邑。这就是变法反复出现的背景。

但“国家能力”不是中性的好词。把户籍、土地、军功和法律拢到中心，能提高战争效率，也会把个人、家族和地方社会卷进更严密的动员。秦最终赢得统一，正是这场竞争的结果；秦朝迅速崩解，也与这种动员如何被用于和平时期有关。
